% 论文摘要（对应 html/src/main.md 标题与目录之间的摘要切片）
\emph{本文运用随机过程的白噪声激励成形表示（原称 Bode-Shannon representation）以及动态系统分析的“状态转移”方法，重新审视经典的滤波与预测问题。所得新结果如下：}

\emph{(1) 问题的数学表述与求解方法无需修改即可同时适用于平稳与非平稳统计特性，以及增长记忆与无限记忆滤波器。}

\emph{(2) 导出了最优估计误差协方差矩阵所满足的一条非线性差分（或微分）方程。由该方程的解无需进一步计算即可直接获得线性最优滤波器的差分（或微分）方程系数。}

\emph{(3) 证明了滤波问题是无噪声调节器问题的对偶问题。}

\emph{本文将所建立的新方法应用于两个熟知问题，印证并推广了已有结果。}

\emph{全文论述基本自包含并从第一性原理出发；随机过程理论的基本概念在\hyperref[sec:appendix]{附录}中予以回顾。}
